\scanpage{25}
d'écrire que $u_K$ commute à $\unsure{\mathfrak{g}|\mathcal{G}}_K$, or par hyp. il commute aux $\omega(X)\in\mathfrak{g}_K$ $(X\in\operatorname{Dér}(\unsure{H|\Pi|K}/k_0)\unsure{]|)}$, Soit $\unsure{\mathfrak{u}|\mathfrak{a}|U}\subset\mathfrak{g}$ le plus petit \unsure{sous-espace} vectoriel de $\mathfrak{g}$ tel que $\mathfrak{u}_K$ contienne les $\omega(X)$. Alors il est clair (alg. linéaire) que \deleted{$\unsure{\mathfrak{u}|u_K}$} $u$ commute à $\mathfrak{u}$, donc il suffit que $\mathfrak{u}=\mathfrak{g}$ pour pouvoir conclure (N.B. Il suffirait que l'\unsure{algèbre|Algèbre|algébre} de Lie engendrée par $\mathfrak{u}$ dans $\mathfrak{g}$ soit $\mathfrak{g}$, mais en fait on se convainc facilement que $\mathfrak{u}$ est déjà une $\unsure{k|K}$-algèbre de \unsure{Lie|bie}....). On peut $\unsure{\overline{\mathrm m}|\overline{\mathrm{m\!m}}}$ remplacer $\mathfrak{u}$ par l'algèbre de Lie \unsure{algébrique} qu'elle engendre
% archive sheet number: 24
